\section{引言}

在 943 个 BioASQ 留出测试问题中，强 Hybrid RRF 检索器有超过 6\% 的问题未能在 top 10 中放置任何金标准证据段落——尽管这些证据存在于 top 100 候选池中。这不是检索失败：金标准段落已经进入候选池。这是排序失败：依赖成对相似度的评分函数将医学读者一眼就能认出的相关段落打散在 top-10 截止线之下。

生物医学证据由问题、段落、MeSH 概念、生物医学实体、文档和知识库关系之间的高阶关联共同组织，而成对比较系统性地遗漏了这些关联。一个关于甲状腺激素类似物的问题和一个讨论 DITPA 的段落之间没有任何表层词汇重叠，但临床医生通过 MeSH 层级可以瞬间建立联系：二者都映射到甲状腺相关的树节点。常规词法和稠密检索仅按成对相似度排序，自然丢失了这一连接。

我们提出 KCH-MedRank，一种将这些高阶关系嵌入为可解释排序特征的知识约束学习排序框架。该方法通过 BM25、稠密检索和倒数排名融合保持候选生成的可复现性~\citep{robertson2009bm25,cormack2009rrf}，随后在每个问题的局部候选邻域中融入检索分数、生物医学语义特征、MeSH 层级感知特征、实体与关系特征，以及可选的超图结构特征。局部超图是捕获 n 元概念交互的一种架构方案；扁平知识特征模型和成对图分解作为受控基线。

核心贡献是一个聚焦的证据重排序框架，附带事后特征归因分析。KCH-MedRank 将这些特征输入查询分组 LambdaMART 排序器产生最终证据排序，排序器的特征重要性分数暴露了哪些信号——MeSH 层级路径、共享实体簇、扩散中心性——驱动了每个排序决策。

在 BioASQ 证据检索基准上，KCH-MedRank 在同一增强候选池上相对于真实 MedCPT Cross-Encoder 基线显著提升了 Recall@10（$+0.0157$，$p=0.0076$）。结构消融表明扁平生物医学知识特征已经解释了大部分收益，而将医学知识约束加入局部超图是必要的：没有约束时，原始超图拓扑引入噪声并降低排序质量；加入约束后，退化被逆转，top-10 指标达到统计显著（$p < 0.01$）。超图在 $k{=}5$ 上提供额外早期排序收益，与深层排序位置的扁平知识信号互补。PubMedQA 和 BEIR NFCorpus 辅助诊断确认了证据覆盖改善和重排序稳健性。

本文贡献为：(1)~一个知识约束学习排序框架，融合检索、生物医学语义、MeSH 层级、实体、关系及可选超图特征用于证据重排序；(2)~一个结构消融框架，分离扁平生物医学知识、图拓扑和医学知识约束的效应，揭示后者是图结构提供有用信号而非噪声的必要条件；(3)~一个可解释性分析，将 648 个救回金标准段落追溯到三种知识约束机制，其中 MeSH 层级路径占主导（75.9\%）。
